\section{An irreducible cubic seed survives finite abelian base change}
\label{sec:cubic-seed-abelian-basechange}

\begin{proposition}[AC-G63-selftwist: no character self-twist and abelian-host irreducibility]
\label{AC-G63-selftwist}
Let $K$ be a number field, let $\ell\ne3$, and let
$\rho:G_K\longrightarrow\operatorname{GL}_3(E)$ be an absolutely
irreducible representation over a finite extension $E/\mathbb Q_\ell$.
Suppose that it admits a stable integral lattice whose residual
semisimplification, over an algebraic closure of the residue field, is
\[
 \bar\rho^{\mathrm{ss}}
       =\bar\chi^{\,j}\bar\gamma\oplus1
                         \oplus\bar\chi^{\,j}\bar\gamma^{-1},
 \qquad j\in\{1,2\},                                      \tag{G63.1}
\]
where $\bar\chi$ has exact order $3$ and the three displayed characters
are pairwise distinct.  Then there is no nontrivial character
$\delta:G_K\to\overline{\mathbb Q}_\ell^{\times}$ such that
$\rho\otimes\delta\simeq\rho$.  Consequently, for every finite
\emph{abelian} extension $L/K$, the restriction $\rho|_{G_L}$ remains
absolutely irreducible.

For the actual cubic residual datum, the proposition applies to a seed
over $K=\mathbb Q(\omega)$ and any of the finite abelian composita
$F=K E_H$ once that seed exists and has (G63.1).  It makes no assertion
that a seed exists, that its automorphic base change has the required
level, or that restriction to an arbitrary finite extension is
irreducible.
\end{proposition}
\begin{proof}
\proofstep{The determinant reduces every possible self-twist to order three.}
If $\rho\otimes\delta\simeq\rho$, taking determinants gives
$\delta^3=1$.  Thus any nontrivial $\delta$ has exact order $3$, takes
values in a finite coefficient extension, and is integral.  Since
$\ell\ne3$, reduction is injective on the third roots of unity:
if a nontrivial third root $z$ reduced to $1$, the relation
$z^2+z+1=0$ would reduce to $3=0$.  It follows that
$\bar\delta$ also has exact order $3$.

\proofstep{The residual character set excludes that possibility.}
An isomorphism with the twist gives equality of traces, and hence
\[
 \bar\chi^{\,j}\bar\gamma+1+
       \bar\chi^{\,j}\bar\gamma^{-1}
 =\bar\delta\bar\chi^{\,j}\bar\gamma+\bar\delta+
       \bar\delta\bar\chi^{\,j}\bar\gamma^{-1}.             \tag{G63.2}
\]
Distinct characters of a group with values in a field are linearly
independent as functions.  Indeed, in a relation with the smallest
possible nonzero number of terms, replace its argument $g$ by $hg$
and subtract the original relation multiplied by one character's
value at $h$.  Choosing $h$ on which a second character differs
produces a shorter nonzero relation, a contradiction.

Both sides of (G63.2) consist of three distinct characters, each with
coefficient $1$.  Character independence therefore says that the set
\[
 S=\{\bar\chi^{\,j}\bar\gamma,1,
                                  \bar\chi^{\,j}\bar\gamma^{-1}\}
\]
is stable under multiplication by $\bar\delta$.  This action has no
fixed element: $\bar\delta a=a$ for a character $a$ would imply
$\bar\delta=1$.  Its orbits have length $3$.  Since $|S|=3$ and
$1\in S$, it follows that
$S=\{1,\bar\delta,\bar\delta^2\}$, whose product is $1$.
On the other hand the product of the original three elements is
$\bar\chi^{\,2j}\ne1$.  This contradiction proves the assertion
about self-twists.  The argument also works in residue characteristic
$2$: the coefficient of any character occurring on just one side of
(G63.2) is $1$ or $-1$, and is nonzero.

\proofstep{Finite abelian restriction is semisimple with scalar commutant.}
Extend coefficients to an algebraically closed field of characteristic
zero, write $V$ for the representation space, and put
$H=G_L$ and $D=G_K/H$.  The restriction to $H$ is semisimple.  To see
this directly, choose a nonzero simple $H$-submodule of $V$ of minimal
positive dimension.  Its finitely many conjugates under representatives
of $D$ are again simple $H$-submodules, and their sum is a nonzero
$G_K$-submodule, hence equals $V$.  A sum of simple submodules is
semisimple: successively add a summand only if its intersection with
the sum already chosen is zero.

Conjugation by $\rho(G_K)$ defines an action of $D$ on
$\operatorname{End}_H(V)$.  It is well defined because elements of
$H$ centralize every member of this endomorphism algebra.  Since $D$
is finite abelian in characteristic zero, this action is the direct
sum of its character eigenspaces.  A member $T$ of the eigenspace for
$\eta:D\to\overline{\mathbb Q}_\ell^\times$ satisfies
\[
 \rho(g)T\rho(g)^{-1}=\eta(g)T.
\]
Thus $T$ is an intertwiner between $\rho$ and a character twist of
$\rho$.  A nonzero such intertwiner is an isomorphism by irreducibility.
The first part excludes every $\eta\ne1$, while the $\eta=1$
eigenspace consists of the scalars by Schur's lemma.  Therefore
$\dim\operatorname{End}_H(V)=1$.
If a semisimple representation is a direct sum of distinct simples
with multiplicities $m_i$, its endomorphism algebra has dimension
$\sum_i m_i^2$.  The dimension just obtained forces one simple
summand with multiplicity one.  This proves absolute irreducibility
over $L$.
\end{proof}

\paragraph{Actual source interface and geometric dictionary.}
The literal seed condition still includes irreducibility:
\path{DenseThresholdCases/IsSeed.lean},
\texttt{IsSeedAlt.irr}; the propagation input still requires an
irreducible determinant specialization:
\path{SecondThirdPrimeConditionTransferAlt.lean},
\texttt{IdealSpecializationPoint}.
The proposition preserves this representation-theoretic property for
a seed already obtained; it supplies neither of those existence
premises.  For finite abelian fields used here, restriction is the
pullback of the associated lisse system along the finite cover after
removing ramified primes, in the previously audited SGA~1,
Expos\'e~V, \S\S4 and~7 framework.  No new archival locator or
general nonabelian base-change claim is used.
